\section{Remarks and open problems}\label{sec:remarks}

% ledger: P1-P8, H12, N3, N4, N5, N6, Q5; source COMMUNITY_PAPER_PROFILE.md §10; 2303.10798 L253, L262; 2409.15880 L49, L236; 2305.17743 L115, L118
We close with questions, in the hat paper's manner, and we phrase each as a question.

\begin{enumerate}[label=\textup{(\arabic*)}, leftmargin=*]
  \item \emph{Symmetry.} The fixed-point construction (\cref{prop:seedsym}) yields tilings with
    symmetry groups of orders $1$ and $8$. Is the bound $|\Sym(T)|\le24$ attained, and which
    other finite symmetry groups occur? Question~32 of Coulbois, Gajardo, Guillon and Lutfalla
    asks for a geometric monotile whose tilings all have trivial stabilizer \cite{CGGL24};
    $\solid$ is not one (\cref{sec:aperiodicity}).
  \item \emph{Fault surfaces and the tiling space.} The ancestors of one fixed tile can fail to
    cover space (\cref{sec:hierarchy}). Which non-exhausting hierarchies occur, and does every
    tiling lie in the substitution hull? The hat paper asks the latter of hats \cite{SMKGS24}; for
    $\solid$ the question is also where a hull-level version of \cref{cor:limitperiodic}
    (\cref{sec:limitperiodic}) would
    have to be settled.
  \item \emph{Convexity.} Is there a convex strongly aperiodic monotile in $\R^3$? With
    reflections forbidden, the biprism is convex and weakly but not strongly aperiodic; with
    reflections allowed it admits periodic tilings; $\solid$ is strongly aperiodic and not convex;
    and for translations alone convex tiles cannot be aperiodic \cite{CGGL24}.
  \item \emph{Achirality.} Is there a strongly aperiodic monotile of $\R^3$ with a mirror symmetry,
    or one whose tilings mix both handednesses? The Spectre paper poses the einstein problem for
    each isometry group $G$ \cite{SMKGS24b}; $\solid$ answers it for the full isometry group of
    $\R^3$ and, being strictly chiral, for the group of proper motions at the same time.
  \item \emph{Simplicity.} What is the smallest number of faces, or of vertices, of a strongly
    aperiodic monotile in $\R^3$? $\solid$ has $4{,}272$ triangles because each of its $192$ features
    is a pyramid; within the twelve-parameter family of \cref{sec:solid} the proof does not
    change, and we do not know how far the architecture can be simplified.
  \item \emph{Heesch and isohedral numbers; decidability.} What Heesch numbers and isohedral
    numbers occur for polyhedra in $\R^3$, and is it decidable whether a given polyhedron tiles
    $\R^3$? We prove nothing in this direction.
  \item \emph{A shorter proof.} Simplified proofs of the hat's aperiodicity appeared within
    months \cite{AA25}. The universality half of our proof, \cref{sec:companions,sec:registration},
    is the part we would most like to see shortened; the finite half is a census that a different
    argument might avoid.
  \item \emph{The mechanism elsewhere.} The construction puts small features on a periodic
    carrier so that the features force a registered lattice substitution whose parent language is
    its own fine language. Does the same design yield einsteins in other dimensions, on other
    lattices, or in other crystallographic settings, and is there a version in the plane with a
    different carrier?
\end{enumerate}

% ledger: O1, N4; source 2303.10798 L262 (deflationary close)
The features of $\solid$ have height unit $1/10000$ and maximum height $12/10000$, measured in
carrier-cell units. Everything the theorem says about tilings by $\solid$, that they exist, that
they are hierarchical, that they have no period and no symmetry of infinite order, and that they
never mix a tile with its mirror image, is decided by those pyramids. We do not know whether
fewer or larger features, or a fundamentally simpler solid, can enforce the same conclusions;
$\solid$ may be far from the simplest example.
